Jesus sits on the mountain, and his disciples come to him. He begins to teach them about happiness. He speaks of people who mourn, people who hunger for justice, people who show mercy. Some of the people he calls blessed will suffer persecution. He promises them consolation, mercy and the kingdom of heaven. To the clean of heart he promises the sight of God.

On All Saints we celebrate those in whom these promises have been fulfilled. Revelation describes them standing before God and the Lamb. They come from every nation and language. The elder explains that they have passed through affliction and have washed their robes in the Lamb's blood. They have needed cleansing, and they have received it from Christ. The glory we celebrate has come through his saving death.

The second reading describes our present condition with great tenderness. The Father has loved us and made us his children. John then says that what we shall be has not yet appeared. We already belong to God, and we await a likeness to him that we cannot yet fully understand. John connects that likeness with seeing him as he is. The future he promises is communion with God himself.

Then John says that everyone who has this hope purifies himself. Those words place a demand on the children God already loves. We have received a life that must grow. Our desire, our choices and our treatment of other people need to become fit for the communion we hope to share. Jesus' promise to the clean of heart concerns that preparation as well as its fulfillment.

Jesus makes this promise to the heart. Saint Augustine explains that a pure heart seeks God without divided desire. We can intend to do what God asks and at the same time insist on keeping something he asks us to surrender. That division becomes a difficulty in loving him. As we hope to see God, we need to learn to love what he loves and give up what keeps us from him.

Augustine explains this waiting with the image of a container being stretched to hold a large gift. The recipient has to become capable of receiving what is promised. Our capacity to know and love God can grow. We discover that things we have treated as our final happiness cannot give us all we seek. We learn to desire God more steadily, including when his promise seems far beyond what we can presently experience.

Augustine gives a second image. A vessel containing vinegar needs to be emptied and cleaned before it can receive honey. Some of our desires need that kind of change. A heart intent on being admired will find it difficult to do good without recognition. A heart attached to resentment will resist the mercy that another person needs. Growing toward God involves giving up desires that keep directing us away from him.

This can be painful because an attachment has become something we want to defend. We can know that mercy is good and still want someone else to remain in our debt. We can ask God for peace and still want the satisfaction of having the last word. The hope of seeing God reaches into those desires. We are preparing to love him with a heart that willingly receives what he loves.

Augustine also explains how we can do what John asks. God purifies us, and we consent to his action. Our will really has something to do. We turn toward the help being offered, acknowledge the sin we have defended, and begin to act differently. Asking for help already expresses that willingness. The person who asks is depending on God and undertaking the work with him.

We do this together in the opening prayer of today's Mass. We ask God for mercy with the saints' intercession. We rejoice in what he has accomplished in them and ask him to help us. Their holiness gives us reason to trust his generosity. Their intercession also means that those whose charity has reached fulfillment continue to care about those who still need help.

The psalm describes the people who seek God's face by their hands and their hearts. Our deeds and our desires both matter. A sincere wish to know God has consequences for what we do to our neighbor. In the Gospel Jesus names mercy before the promise of vision. The people preparing to see God learn to show mercy while they themselves need it.

Augustine makes that obligation very plain. You may need God's help and still be able to help someone weaker than you. Your own need gives you a reason to be merciful. You know something of what it means to depend on another's kindness. There is a good you can do for that person even while much in your own life remains unfinished.

You can put that mercy into practice by giving needed help to someone whose faults have made you reluctant to help. Give the help you can honestly give, without making the person endure a rehearsal of those faults. That costs something when resentment has become familiar. It is also a way of consenting to God's work in your own heart. The person you help receives a real good, and you begin to relinquish a desire that has kept you from charity.

The promise of seeing God gives us a reason to persevere when that change is slow. In Revelation, the multitude has come through affliction. Its members praise God and the Lamb for their salvation. Bede explains that cleansing in the Lamb's blood includes the whole Church, beyond those who died as martyrs. Christ's saving work reaches people whose fidelity has taken different forms. We can depend on him in the work of purification still before us.

Jesus likewise promises blessedness amid circumstances that remain difficult. Mourning awaits consolation. Hunger for justice awaits satisfaction. Mercy is practiced by people who themselves need mercy. A disciple may begin to live these words and still have much to suffer and much to learn. The promised end gives that unfinished life a direction. We seek the sight of the God who has already made us his children.

At this Mass we receive Communion, and afterward the Church still asks for charity to be brought to completion. The prayer looks toward the heavenly banquet. We receive Christ now and need his grace to continue growing in love. The saints' fulfillment shows us what we hope to share. Their concern for us belongs to the same charity for which we ask.

You come to this Eucharist as people whom the Father loves. You can acknowledge your need for cleansing in the presence of the Lamb who saves. He has brought the saints through their affliction into his joy. As you receive him, entrust your unfinished life to his help and consent to the love he is teaching you. We hope to stand with the saints before God, able to love him and one another with the charity he brings to completion.
